%% ============================================================
%%  The Course List of Facts.
%%  Back-matter reference page. The facts use stable keys and
%%  auto-generated F## identifiers supplied by \coursefact.
%% ============================================================
\documentclass[main.tex]{subfiles}
\begin{document}

\clearpage
\phantomsection\addcontentsline{toc}{chapter}{The Course List of Facts}
\markboth{The Course List of Facts}{The Course List of Facts}
\thispagestyle{plain}
{\LARGE\bfseries\hyphenpenalty=10000\exhyphenpenalty=10000
 The Course List of Facts\par}\nobreak\smallskip
\noindent\rule{\linewidth}{0.4pt}\par\medskip
\noindent This is the student's ledger of quotable results: the rules and definitions used
without reproof in solutions and derivations, seeded from Lecture~1 and Assignment~1 and
growing week by week. Each fact has a stable key and an automatically generated identifier,
and each carries a pointer to where it lives in these notes. ``Quotable'' means you may use it
without reproving it, but you should still know \emph{why} it holds.

\par\medskip\noindent{\large\bfseries Probability Rules}\par\nobreak\smallskip
\begin{itemize}[leftmargin=1.4em,itemsep=2pt,topsep=2pt]
  \coursefact{lin-exp}
    {Expectation is linear: $\E\{aX+bY+c\}=a\,\E\{X\}+b\,\E\{Y\}+c$, regardless of any
     dependence between $X$ and $Y$.}
    {\xrefL{1.4.1}; Problem 1.13}
  \coursefact{lotus}
    {The law of the unconscious statistician: $\E\{g(Y)\}=\sum_y g(y)\,p(y)$ or
     $\int g(y)f(y)\,dy$ --- compute against the distribution of $Y$ itself.}
    {\xrefL{1.4.2}}
  \coursefact{var-def}
    {$\Var(Y)=\E\{(Y-\E\{Y\})^2\}=\E\{Y^2\}-(\E\{Y\})^2$; for a Bernoulli variable,
     $\Var(Y)=\theta(1-\theta)$.}
    {Problem 1.9; \xrefL{1.A.2}}
  \coursefact{var-affine}
    {$\Var(aY+c)=a^2\,\Var(Y)$: constants shift, scales square.}
    {Problem 1.14; \xrefL{1.A.3}}
  \coursefact{var-sum}
    {$\Var(U+V)=\Var(U)+\Var(V)+2\,\Cov(U,V)$; for independent $U,V$ the covariance term
     vanishes and variances add.}
    {Problem 1.14; \xrefL{1.A.2}}
  \coursefact{cov-bilinear}
    {Covariance is bilinear and blind to constants:
     $\Cov(aU+b,\,cV+d)=ac\,\Cov(U,V)$, and covariances of sums expand term by term.}
    {Problem 1.15}
  \coursefact{cov-self}
    {$\Cov(Y,Y)=\Var(Y)$.}
    {Problem 1.15}
  \coursefact{indep-factor}
    {Independence means the joint distribution factors into the marginals; in particular
     $\E\{XY\}=\E\{X\}\,\E\{Y\}$ for independent $X,Y$.}
    {Problem 1.11}
  \coursefact{indep-cov}
    {Independence implies $\Cov(X,Y)=0$; the converse is false (e.g.\ $X$ uniform on
     $\{-1,0,1\}$ and $Y=X^2$).}
    {Problem 1.12}
  \coursefact{iterated-exp}
    {Iterated expectations: $\E\{\E\{Y\mid X\}\}=\E\{Y\}$.}
    {\xrefL{1.4.6}}
  \coursefact{beta-facts}
    {For the Beta$(\alpha,\beta)$ distribution on $[0,1]$: mean $\alpha/(\alpha+\beta)$ and
     variance $\alpha\beta/\big((\alpha+\beta)^2(\alpha+\beta+1)\big)$; Beta$(1,1)$ is the
     uniform distribution.}
    {\xrefL{1.8.5}}
\end{itemize}

\par\medskip\noindent{\large\bfseries Decision Theory}\par\nobreak\smallskip
\begin{itemize}[leftmargin=1.4em,itemsep=2pt,topsep=2pt]
  \coursefact{model-def}
    {The statistical model is the set of all possible distributions of the data, usually
     parameterized as $\{p^{\theta}$ or $f^{\theta}:\theta\in\Theta\}$.}
    {\xrefL{1.3.1}, \xrefL{1.3.2}}
  \coursefact{risk-def}
    {The risk of a decision rule $\delta$ is the expected loss as a function of the parameter,
     $R_{\delta}(\theta)=\E^{\theta}\{L(\delta(Y),\theta)\}$.}
    {\xrefL{1.5.4}}
  \coursefact{admissible}
    {A method is admissible if no other method does at least as well for every parameter value
     and better for some.}
    {\xrefL{1.5.6}}
  \coursefact{bayes-risk}
    {The Bayes risk of $\delta$ under a prior $\pi$ is the expected risk,
     $\E\{R_{\delta}(\theta)\}=\int_{\Theta}\pi(\theta)\,R_{\delta}(\theta)\,d\theta$; a Bayes
     rule minimizes it.}
    {\xrefL{1.6.2}, \xrefL{1.6.3}}
  \coursefact{admissible-bayes}
    {As stated in lecture: ``(almost) every admissible rule is a Bayes rule and every Bayes
     rule is admissible'' --- heuristic form; precise versions need side conditions.}
    {\xrefL{1.6.5}}
  \coursefact{bias-variance}
    {Under squared-error loss, risk decomposes as
     $R_{\delta}(\theta)=\Var^{\theta}(\delta)+\text{bias}(\theta)^2$.}
    {\xrefL{1.A.1}}
  \coursefact{beta-posterior}
    {Binomial data with a Beta$(1,1)$ prior give a Beta$(1+x,\,1+n-x)$ posterior, whose mean
     $(1+x)/(n+2)$ is the Bayes rule for squared-error loss.}
    {\xrefL{1.A.4}}
\end{itemize}

\end{document}
